% Generated draft: requires mathematical and source review.
\documentclass[11pt]{article}
\usepackage[margin=25mm]{geometry}
\usepackage{fontspec}
\setmainfont{Latin Modern Roman}
\usepackage{amsmath,amssymb,mathtools,unicode-math}
\setmathfont{Latin Modern Math}
\usepackage{enumitem,needspace,xurl,hyperref}
\setlength{\parindent}{0pt}
\setlength{\parskip}{5pt}
\setlength{\emergencystretch}{3em}
\newcommand{\N}{\mathbb{N}}
\newcommand{\Z}{\mathbb{Z}}
\newcommand{\Q}{\mathbb{Q}}
\newcommand{\R}{\mathbb{R}}
\newcommand{\C}{\mathbb{C}}
\newcommand{\F}{\mathbb{F}}
\newcommand{\A}{\mathbb{A}}
\newcommand{\PP}{\mathbb{P}}
\newcommand{\E}{\mathbb{E}}
\newcommand{\eps}{\varepsilon}
\newcommand{\dd}{\,\mathrm{d}}
\begin{document}
\section*{UnsolvedMath ID: 701021}
\textbf{Source problem code:} catalogue entry
\subsection{Definitions and mathematical statement}
A subgroup \(H\leq G\) is subnormal in \(G\) if there are an integer \(n\geq 0\) and subgroups
\[
H=H_0\trianglelefteq H_1\trianglelefteq\cdots\trianglelefteq H_n=G.
\]
A group \(G\) is a \(T\)-group if every subnormal subgroup of \(G\) is normal in \(G\). It is an \(\overline{T}\)-group if every subgroup \(K\leq G\), including \(G\) itself, is a \(T\)-group; thus, for all \(H\leq K\leq G\), subnormality of \(H\) in \(K\) implies \(H\trianglelefteq K\).

A group \(G\) is locally graded if every nontrivial finitely generated subgroup \(K\leq G\) has a proper subgroup \(L<K\) of finite index, so that \(1<[K:L]<\infty\). A group is periodic if every element has finite order; hence \(G\) is non-periodic if it contains an element of infinite order.

The question is whether, for every group \(G\),
\[
\bigl(G\text{ is non-periodic, locally graded, and an }\overline{T}\text{-group}\bigr)
\quad\Longrightarrow\quad
\forall x,y\in G,\;xy=yx.
\]
The supplied registry labels this question open; that status is supplied source metadata and has not been independently verified.
\subsection{Short English statement}
Must every locally graded group containing an element of infinite order be abelian if every subgroup is a \(T\)-group, meaning that its subnormal subgroups are normal in it?
\subsection{Sources}
\begin{itemize}
\item Local UnsolvedMath registry, supplied complete source record: permanent ID \(701021\), exact source problem code \texttt{catalogue entry}, title \emph{Kourovka Notebook Problem 14.36}. The supplied status is \texttt{open}.
\item F. de Giovanni, proposer, 1999. The supplied attribution identifies \emph{The Kourovka Notebook (Unsolved Problems in Group Theory)}, 14th issue (1999), original problem number 14.36, and the 21st issue edition (2026), \texttt{arXiv:1401.0300v48}. The supplied collection description names E. I. Khukhro and V. D. Mazurov as editors. These bibliographic details are supplied attribution and have not been independently verified.
\end{itemize}
\end{document}
